%=============================================================================!
% 1.7  CIRCULAR MOTION                                                        !
%=============================================================================!
\section{Circular motion}
\label{sec:mo-circle}

A stone whirled on a string at a steady rate moves around a circle of
radius $R$, centred on the origin, at constant speed. Its angle from the
$x$ axis grows at a steady rate $\omega$, the angular speed in radians
per second, and we start the clock as it crosses the positive $x$ axis,
so that the angle is $\omega t$ at time $t$. By the definition of sine and cosine in
\cref{sec:mo-trig}, scaled to a circle of radius $R$, the position of the
stone is
\begin{equation*}
   \vb{r}(t) = R\,(\cos\omega t,\ \sin\omega t) ,
\end{equation*}
and it completes one revolution when $\omega t$ has grown by $2\pi$, in
the period $2\pi/\omega$.

\subsection*{Velocity and acceleration}

Each component is a sinusoid, and the chain rule with the derivatives of
\cref{eq:mo-trig-derivs} gives
\begin{equation*}
   \frac{\dd}{\dd t}\cos\omega t = -\omega\sin\omega t, \qquad
   \frac{\dd}{\dd t}\sin\omega t = \omega\cos\omega t .
\end{equation*}
Differentiating the position once and then a second time therefore gives
\begin{equation}
   \vb{v}(t) = R\omega\,(-\sin\omega t,\ \cos\omega t), \qquad
   \vb{a}(t) = -R\omega^2\,(\cos\omega t,\ \sin\omega t) = -\omega^2\vb{r}(t).
   \label{eq:mo-circle}
\end{equation}
The scalar product of the velocity with the position is
\begin{equation*}
   \vb{v}\cdot\vb{r} = R^2\omega\,\bigl(-\sin\omega t\cos\omega t
   + \cos\omega t\sin\omega t\bigr) = 0 ,
\end{equation*}
so the velocity is perpendicular to the radius and hence tangent to the
circle, as \cref{sec:mo-plane} requires. Its magnitude is
\begin{equation*}
   v = R\omega\,\bigl(\sin^2\omega t + \cos^2\omega t\bigr)^{1/2} = R\omega ,
\end{equation*}
which does not change in time. The radian, being a ratio of lengths,
drops out of the units, so $R\omega$ is in metres per second. The
acceleration $-\omega^2\vb{r}$ points
from the stone towards the centre, opposite to the position vector, with
magnitude $\omega^2 R$. Since $\omega = v/R$, this magnitude can also be
written $\omega^2 R = v^2/R$ (\cref{fig:mo-plane}b). A stone on a string
of \SI{0.5}{\metre} going round twice a second, for example, has $\omega
= 2\times2\pi = \SI{12.57}{\radian\per\second}$, a speed of
$\SI{6.28}{\metre\per\second}$ and an acceleration of
$\SI{79}{\metre\per\second\squared}$ towards the centre, about eight
times $g$.

\subsection*{Constant speed with non-zero acceleration}

Uniform circular motion therefore has a constant speed but a non-zero
acceleration at every instant, because acceleration measures any change of
velocity, including a change of direction alone. The relation between
speed and acceleration found in \cref{sec:mo-acceleration} extends to
vectors. Writing $v^2 = \vb{v}\cdot\vb{v} = v_x^2 + v_y^2 + v_z^2$ and
applying that result to each component,
\begin{equation*}
   \frac{\dd}{\dd t}\bigl(v_x^2 + v_y^2 + v_z^2\bigr)
   = 2v_x a_x + 2v_y a_y + 2v_z a_z = 2\,\vb{v}\cdot\vb{a} .
\end{equation*}
The speed is therefore constant over an interval exactly when
$\vb{v}\cdot\vb{a} = 0$ at every instant of that interval, which is the
case on the circle. On the parabola of \cref{sec:mo-plane}, by contrast,
$\vb{v}\cdot\vb{a} = v_0\cos\alpha\times0 + (v_0\sin\alpha - gt)(-g) =
-g\,v_y$, where $v_y = v_0\sin\alpha - gt$ is the vertical velocity. This
is negative while the ball rises ($v_y > 0$) and positive while it falls,
so the ball slows as it rises and speeds up as it falls; only at the top
is it zero.

\subsection*{The projection onto a line}

The $x$ coordinate of the stone, the projection of its motion onto the $x$
axis, moves back and forth along a line as $x(t) = R\cos\omega t$, a
sinusoid of amplitude $R$ and angular frequency $\omega$. The $x$
component of \cref{eq:mo-circle} gives this projection the acceleration
$-R\omega^2\cos\omega t$, which is $-\omega^2$ times $x(t)$ itself:
\begin{equation*}
   \frac{\dd^2 x}{\dd t^2} = -\omega^2 x ,
\end{equation*}
the property of every sinusoid found in \cref{sec:mo-trig}. The
acceleration of the projection is proportional to its distance from the
centre and directed back towards it, so it is greatest at the ends of the
swing and zero as it passes through the middle. \Cref{sec:ne-spring}
shows that a block attached to a spring obeys this equation, and
Chapter~4 that the small vibrations of atoms obey it to a good
approximation.

\notebook{01}{circular motion and its projection}
